% Einheit 2 von 3 – Kettenregel (bank/ableitungsregeln/e2.jsonl, zusammenbau v0.1)
%% TODO Zeitmarke Einheit 2: Marken-Zeile nicht lesbar
%% TODO Zweigzeile Teil 1: was hier gelernt wird (Satz in Schülersprache aus dem Einheitstitel) – Einheit 2
\einheitenkopf[e2]{Einheit 2 von 3 · Kettenregel}
\zweigzeile{keine Prüfungsaufgabe}
\setcounter{aufgabe}{13}

%% TODO Titel: Ich-kann-Satz fehlt in der Bank (Platzhalter: Kettenname) – Nr. 14
%% TODO Anweisung: ein Satz über der ersten Teilaufgabe fehlt in der Bank – Nr. 14
\begin{aufgabe}{Kettenregel}
%% TODO \rechenplatz{n}: Schrittzahl des Grundfalls fehlt in der Bank (nur bei fester Schreibform, 2.3 b)
\begin{teile}
\teil $f(x) = e^{7x}$ – welcher Teil ist die innere Funktion? \\ \kreuz{der Exponent $7x$} \\ \kreuz{die Basis $e$} \\ \kreuz{die Zahl $7$}
\teil $f(x) = e^{5x}$ – $f'(x)$? $f'(x) =$ \leerfeld
\teil $f(x) = (3x + 1)^4$ – $f'(x)$? $f'(x) =$ \leerfeld
\teil $f(x) = e^{2x^2}$ – $f'(x)$? $f'(x) =$ \leerfeld
\teil $f(x) = e^{4x}$ – Term der fünften Ableitung $f^{(5)}$? Um wie viel muss man den Graphen von $f$ in $x$-Richtung verschieben, damit der Graph von $f^{(5)}$ entsteht? \hfill (Abitur 2024 LK) \\ $f^{(5)}(x) =$ \leerfeld , Verschiebung: \leerfeld
\teil Die Abbildung zeigt die Graphen der ganzrationalen Funktionen $f$ und $g$; $g$ hat den Tiefpunkt $T$. Für $h(x) = g(f(x))$: Gleichung der Tangente an den Graphen von $h$ im Punkt $(2|h(2))$? \hfill (Abitur 2020 LK) \\

\begin{ksys}[xmin=-3,xmax=4,ymin=-3,ymax=5,ablesen] \funktion{0.5*\x^2-1}{f} \parabel{1}{1}{-2}{g} \tiefpunkt{1}{-2}{T} \end{ksys}
Tangente: \leerfeld
\end{teile}
\end{aufgabe}

%% TODO Titel: Ich-kann-Satz fehlt in der Bank (Platzhalter: Kettenname) – Nr. 15
%% TODO Anweisung: ein Satz über der ersten Teilaufgabe fehlt in der Bank – Nr. 15
\begin{aufgabe}{Term einer höheren Ableitung einer Sinusfunktion über die Periodizität der Ableitungen angeben}
\begin{teile}
\teil $f(x) = \mathrm{sin}(2x)$ – Term der fünften Ableitung $f^{(5)}$? \hfill (Abitur 2017 LK) Es gilt $\mathrm{sin}' = \mathrm{cos}$ und $\mathrm{cos}' = -\mathrm{sin}$. $f^{(5)}(x) =$ \leerfeld
\end{teile}
\end{aufgabe}

%% TODO Titel: Ich-kann-Satz fehlt in der Bank (Platzhalter: Kettenname) – Nr. 16
%% TODO Anweisung: ein Satz über der ersten Teilaufgabe fehlt in der Bank – Nr. 16
\begin{aufgabe}{Kettenregel – Fehler finden, Begründen}
\begin{teile}
\teil Ich finde den Fehler: Zu $f(x) = e^{-6x}$ schreibt Jan $f'(x) = e^{-6x}$.
\teil Begründe, warum $e^{4x}$ die Ableitung $4 \cdot e^{4x}$ hat und nicht $e^{4x}$.
\end{teile}
\end{aufgabe}
